% =====================================================================
%  CARNET D'ENTRAINEMENT — FICHE 6 — THÈME 3
%  Niveau MOYEN — Exercices
% =====================================================================
\documentclass[11pt,a4paper]{article}
\usepackage[utf8]{inputenc}
\usepackage[T1]{fontenc}
\usepackage{lmodern}
\usepackage[french]{babel}
\usepackage{amsmath,amssymb,mathtools}
\usepackage{icomma}
\usepackage[version=4]{mhchem}
\usepackage{siunitx}
\sisetup{output-decimal-marker={,}}
\usepackage{xcolor}
\usepackage{colortbl}
\usepackage{tikz}
\usepackage[most]{tcolorbox}
\usepackage{enumitem}
\usepackage{array}
\usepackage{booktabs}
\usepackage{fancyhdr}
\usepackage[hidelinks]{hyperref}
\usetikzlibrary{arrows.meta,positioning,calc,shapes.geometric}
\usepackage[a4paper,margin=2.1cm,headheight=26pt,headsep=8pt]{geometry}

\definecolor{primary}{RGB}{146,95,160}
\definecolor{primarydark}{RGB}{92,52,104}
\definecolor{excol}{RGB}{0,135,90}
\definecolor{atomO}{RGB}{220,55,45}
\definecolor{atomN}{RGB}{48,80,200}
\definecolor{atomH}{RGB}{235,235,235}
\definecolor{lonepair}{RGB}{120,94,200}
\definecolor{bondcol}{RGB}{60,60,70}

\tcbset{boxbase/.style={enhanced,breakable,boxrule=0.9pt,arc=2.5pt,
   left=3mm,right=3mm,top=2.2mm,bottom=2.2mm,
   fonttitle=\bfseries,coltitle=white,toptitle=0.6mm,bottomtitle=0.6mm}}
\newtcolorbox[auto counter]{exo}[1][]{boxbase,colback=primary!5,colframe=primary,
   title={\textsc{Exercice}~\thetcbcounter\ifx\relax#1\relax\relax\else\textmd{ —\itshape\ #1}\fi}}

\pagestyle{fancy}\fancyhf{}
\fancyhead[L]{\small\textcolor{primarydark}{\textbf{Fiche 6 · Thème 3} — Charge formelle et choix de structure}}
\fancyhead[R]{\small\textcolor{primarydark}{\textsc{Moyen}}}
\fancyfoot[C]{\small\thepage}
\renewcommand{\headrulewidth}{0.8pt}
\renewcommand{\headrule}{\hbox to\headwidth{\color{primary}\leaders\hrule height \headrulewidth\hfill}}
\setlength{\parindent}{0pt}\setlength{\parskip}{3pt}

\begin{document}
\begin{center}
{\Large\bfseries\color{primarydark}Thème 3 — Niveau Moyen}\\[2pt]
{\small\itshape On calcule la charge formelle de chaque atome d'une structure \emph{donnée}, puis
on vérifie la cohérence. Rappel : une double liaison compte pour 2 liaisons.}
\end{center}
\medskip

\begin{exo}[ion oxonium \ce{H3O+}]
Dans l'ion oxonium, l'oxygène central forme \textbf{3 liaisons \ce{O-H}} et porte \textbf{1 doublet
non liant} ; chaque \ce{H} forme une liaison simple.
\begin{enumerate}[label=\alph*),leftmargin=2em,itemsep=1pt]
  \item Calculer la charge formelle de l'oxygène.
  \item Calculer la charge formelle de chaque hydrogène.
  \item Faire la somme et vérifier qu'elle est cohérente avec la charge de l'ion.
\end{enumerate}
\end{exo}

\begin{exo}[ion carbonate \ce{CO3^2-}]
La structure comporte un carbone central lié à trois oxygènes : \emph{un} \ce{O} doublement lié
(2 doublets libres) et \emph{deux} \ce{O} simplement liés (3 doublets libres chacun).
\begin{enumerate}[label=\alph*),leftmargin=2em,itemsep=1pt]
  \item Calculer la charge formelle du carbone, de l'oxygène doublement lié et d'un oxygène simplement lié.
  \item Vérifier que la somme des charges formelles vaut $-2$.
  \item Les charges négatives sont-elles \emph{bien} placées ? Justifier avec l'électronégativité.
\end{enumerate}
\end{exo}

\begin{exo}[ion ammonium \ce{NH4+}]
L'azote de l'ion ammonium forme \textbf{4 liaisons \ce{N-H}} et ne porte \textbf{aucun doublet libre}.
\begin{enumerate}[label=\alph*),leftmargin=2em,itemsep=1pt]
  \item Calculer la charge formelle de l'azote et celle de chaque hydrogène ; vérifier la somme.
  \item Expliquer pourquoi c'est l'azote « tétravalent » qui porte la charge $+$ de l'ion, alors
        même que c'est lui qui avait fourni le doublet de la liaison dative.
\end{enumerate}
\end{exo}

\begin{exo}[molécule d'ozone \ce{O3}]
L'ozone a la structure \ce{O=O-O} : l'atome d'oxygène \textbf{central} est lié par une
\textbf{double liaison} à un oxygène et par une \textbf{simple liaison} à l'autre, et il porte
\textbf{1 doublet non liant}. L'oxygène \emph{doublement lié} porte 2 doublets libres ; l'oxygène
\emph{simplement lié} porte 3 doublets libres.
\begin{enumerate}[label=\alph*),leftmargin=2em,itemsep=1pt]
  \item Calculer la charge formelle de l'oxygène central. \emph{(Attention : la double liaison
        compte pour 2 liaisons — l'atome central a donc $2+1 = 3$ liaisons.)}
  \item Calculer la charge formelle de l'oxygène doublement lié, puis de l'oxygène simplement lié.
  \item Vérifier que la somme est nulle, cohérente avec une molécule neutre.
\end{enumerate}
\end{exo}

\begin{exo}[comparer deux structures d'un même atome]
Pour une molécule dont l'atome central est un \textbf{soufre} (période 3), deux structures de Lewis
sont envisagées : dans la structure (I), le soufre porte une charge formelle $q_{\mathrm{f}} = +2$ ;
dans la structure (II), il porte $q_{\mathrm{f}} = 0$.
\begin{enumerate}[label=\alph*),leftmargin=2em,itemsep=1pt]
  \item Laquelle des deux structures est la plus représentative ? Énoncer le critère utilisé.
  \item Par quel mécanisme le soufre a-t-il pu faire passer sa charge formelle de $+2$ à $0$ ?
\end{enumerate}
\end{exo}

\end{document}
